\chapter{Globulin Test}
\section*{What Is This Test?} A globulin test estimates blood proteins other than albumin and helps evaluate immune, liver, kidney, and nutritional conditions.\section*{Alternative Names} Serum globulin; total protein and albumin/globulin ratio; A/G ratio.\section*{Principle} Globulin is often calculated from total protein minus albumin; electrophoresis separates protein fractions.\section*{Why This Test Is Ordered} It may investigate recurrent infection, inflammation, liver disease, protein loss, malnutrition, or a suspected monoclonal protein.\section*{Primary vs Secondary Test} It is a screening chemistry result; electrophoresis and immunofixation characterize abnormalities.\section*{How This Test Is Performed} A venous blood sample is analyzed for total protein and albumin, with calculation or direct fraction testing.\section*{What Preparation Is Needed} Fasting depends on the accompanying panel. Report hydration status and medicines.\section*{Understanding Results} High globulin may reflect inflammation or monoclonal protein; low values may reflect protein loss, reduced production, or immune deficiency.\section*{Follow-up Tests} Serum protein electrophoresis, immunofixation, immunoglobulins, liver and kidney tests, or urine protein studies may follow.\section*{References} \begin{enumerate}\item MedlinePlus. Globulin Test.\item College of American Pathologists. Serum protein electrophoresis resources.\end{enumerate}\clearpage\section*{Understanding Results and Follow-up}
The laboratory report should identify the specimen, method, units, reference interval or decision limit, and whether the result is positive, negative, abnormal, or inconclusive. Reference intervals vary with age, sex, pregnancy, specimen type, assay, and laboratory; a result outside a range is not automatically a diagnosis, and a result within a range does not always exclude disease. Discuss unexpected results with the ordering clinician, who can decide whether repeat or confirmatory testing, treatment, or referral is appropriate. Do not change medicines or delay urgent care based on an isolated result.
\section*{Regulatory Status and Authorization Date}This chapter describes a test category, not a specific commercial product. FDA, CDSCO, EU IVDR, and MHRA approval or registration dates therefore cannot be assigned without the assay manufacturer, product name, instrument, and intended use. For laboratory-developed tests, an individual agency approval date may not exist. See the Preface for the interpretation of regulatory status.
\clearpage
